\section{Réactions du Conseil Municipal}

« Ce hameau forestier » c’est l’appellation officielle, avait été édifié, comme les 74
autres hameaux dans des conditions très sommaires, écrit Michel Roux dans son livre
« Les harkis, les oubliés de l’histoire 1954-1991 » P.312 – La Découverte- La qualité médiocre
des matériaux de construction, la précarité, l’insalubrité et l’exiguïté des logements,
les défauts de viabilisation sont le lot de la quasi-totalité de ces regroupements
théoriquement conçus pour s’effacer rapidement de la carte...

A Meyrueis, à Truscas (Hérault) et Mirande (Gers) – comme disent Abi Slimane et Finas
dans le livre -Regroupement et dispersion-) le chauffage est insuffisant et l’humidité pénètre
partout. Avec leur défaut d’isolation les logements sont travaillés de l’intérieur, par
en dessous comme par en dessus ? par une pourriture insistante qui les ronge à un
rythme accéléré. Le chauffage ? quand il parvient (mais à un prix exorbitant) à
réchauffer les habitants, génère une sorte de brouillard, une vapeur « tropicale-
humide » : micro-climat du micro-monde forestier ; le hameau est le cocon morbide
du Français-Musulman ».

Leur séjour à Meyrueis ne laissa que peu de traces dans les archives municipales. Il
faut dire que ces hameaux forestiers ou camps de harkis relevaient directement du
Ministère des Rapatriés et du Ministère de l’Agriculture dont dépend le Service des
Eaux et Forêts. Les communes n’avaient pas leur mot à dire, les ministères
s’entendant directement avec les Eaux et Forêts où beaucoup seront employés. J’ai
noté seulement deux paragraphes dans les délibérations du Conseil Municipal de
Meyrueis :

\begin{quote}
« Le 11 Décembre 1962, le conseil municipal est appelé à donner son avis sur le
transport des élèves de familles harkies qui doivent être installées dans la commune.
Le conseil municipal donne un avis favorable pour le ramassage de ces élèves sous
réserve que cette opération n’occasionne aucun frais pour la commune ».
\end{quote}

Un peu plus tard le 30 Juillet 1963 :

\begin{quote}
« L’implantation d’un camp de Harkis à Roquedols, soit 131 personnes, porte la
population de Meyrueis de 1198 à 1329 habitants ; cette augmentation se traduit,
pour la commune, par des charges nouvelles : admission à l’Assistance médicale,
alimentation, eau potable etc. Pour ces raisons le Conseil estime juste que ces
habitants soient ajoutés à la population légale de la commune avec toutes les
conséquences de droit. »
\end{quote}

\begin{quote}
« Les municipalités, en général, - note toujours Michel Roux (op. cité) - manifestent
d’entrée une réticence non dissimulée à l’implantation de ces Algériens sur leur
territoire communal. Cependant l’éloignement des hameaux et leur encadrement
paramilitaire les rassurent. En outre, à l’assurance de n’avoir pas à subir les charges
imputables à l’arrivée d’une population importante, s’ajoute la promesse qu’elle sera
employée à des travaux d’aménagement dans les territoires communaux. Enfin
nombre de hameaux sont considérés comme temporaires... »
\end{quote}
